\chapter{Criterio mínimo de régimen verificable}

\section{Dominio verificable}
\begin{definition}
Se dice que una configuración o trayectoria entra en dominio verificable cuando sus diferencias, sus pasos y su registro acumulativo pueden ser leídos de forma coherente dentro de una ventana \Win\ y bajo la resolución \eps.
\end{definition}

\begin{observation}
La verificabilidad no equivale todavía a estabilidad fuerte ni a clausura total del régimen. Solo declara que el sistema ha alcanzado un umbral suficiente de legibilidad estructural para ser tratado sin ambigüedad inmediata.
\end{observation}

\section{Criterio mínimo}
\begin{definition}
Una configuración satisface el criterio mínimo de régimen verificable si concurren localmente las siguientes condiciones:
\begin{enumerate}
  \item existe legibilidad de sus componentes;
  \item existe legibilidad de sus pasos estructurales;
  \item existe coherencia mínima dentro del dominio declarado;
  \item existe acumulación mínima no colapsada de forma inmediata.
\end{enumerate}
\end{definition}

\begin{observation}
Estas condiciones no agotan la descripción completa de un régimen. Funcionan como umbral mínimo de admisibilidad estructural.
\end{observation}

\section{Persistencia local}
\begin{definition}
Se llama persistencia local al mantenimiento de legibilidad estructural a través de una sucesión finita de pasos dentro de la misma ventana \Win.
\end{definition}

\begin{observation}
La persistencia local fortalece el paso desde la mera diferencia puntual hacia una primera forma de régimen verificable, pero no implica todavía invariancia de largo alcance.
\end{observation}

\section{Compatibilidad con oposición y proyección}
\begin{definition}
El criterio mínimo de régimen verificable se considera compatible con la involución \Iop\ y con la lectura proyectada cuando ambas operaciones conservan una gramática estructural legible dentro del dominio declarado.
\end{definition}

\begin{observation}
Esta compatibilidad no exige equivalencia total entre estructura primitiva, estructura invertida y proyección. Solo exige que ninguna de estas operaciones destruya de inmediato la verificabilidad local.
\end{observation}

\section{Condición mínima de salida}
\begin{definition}
Se considera que una configuración sale del dominio verificable cuando pierde de forma local alguna de las condiciones mínimas de legibilidad, coherencia o acumulación exigidas por este bloque.
\end{definition}

\begin{observation}
La salida del dominio verificable no clasifica todavía el tipo final de colapso. Solo marca que el régimen ya no puede sostenerse como estructuralmente tratable en el sentido mínimo adoptado.
\end{observation}

\section{Cláusula de alcance}
\begin{observation}
El criterio mínimo de régimen verificable fija un umbral de entrada al tratamiento disciplinado de configuraciones y trayectorias. No equivale aún a una teoría completa de regímenes, estabilidad madura o taxonomía final de comportamientos.
\end{observation}

